\section{Arithmetic terms after the choice of a spectral setting}
Fix the original datum $D$, a setting $X$ and $Y:\mathrm{Full}(X)$.
For words $u,u'$ write $F,F'$ for their flags, and $U_F,U_{F'}$ for
the isotropic spans. All vector spaces and linear maps in this section
are over the original CM field $E$. An adelic term $s$ specifies two
row lattices and a double-unipotent class of rational isometries
$g_s:\Lambda'\longrightarrow\Lambda$. Choose the representative
specified in $s$, and write $G_s=g_sF'$ for its transported flag.
The map $s\mapsto t$ to arithmetic terms remembers the left lattice
and the double equivalence class of the transported frame. The finite
row index remembers the two lattices and a tuple of invertible matrices
$\iota_s$ at the level primes. These two maps forget different data.

% BEGIN ARF1
\paragraph{ARF1. Finiteness of each arithmetic-term fibre.}
Over any fixed arithmetic term $t$, there are finitely many adelic
terms. If the fibre is occupied, its size is at most the order of
the automorphism group of the projected positive definite lattice
of any one right row in the fibre.
\begin{proof}
Fix one lift with return $g$. For another lift with return $g'$,
the definition of the quotient of allowable frames gives two
integral unipotents $n\in N_{\Lambda,F}$, $m\in N_{\Lambda,gF'}$
such that the transported vectors of $g'$ are obtained by applying
$nm$ to those of $g$. Here equality of classes is double equivalence:
the relation is already transitive since conjugating an integral
flag unipotent by a lattice isometry transports its defining flag,
lattice and orthogonal-quotient conditions. Inverses give symmetry.
Equivalently
\[
       g'=n g\mu c,\qquad \mu=g^{-1}mg\in N_{\Lambda',F'},
       \qquad c=(nmg)^{-1}g'.
\]
The isometry $c$ fixes every right isotropic vector $e'_a$ and
the distinguished vector $z'$, and carries the right lattice of $g'$ to that of $g$. The two
right lattices are entries of the same list of representatives, and
$c$ is in the prescribed flag stabilizer. Uniqueness in the row
system therefore makes those two entries equal.

Choose such a $c$ for each lift. On $U_{F'}^\perp/U_{F'}$, it fixes
the negative vector $\bar z'$ and acts on its positive definite
orthogonal complement. It preserves the full projected lattice
in that complement. QPR5 proves finiteness of its isometry group
by taking the positive trace form at all infinite places.
If $c,d$ give the same action on the complement, then $c^{-1}d$
acts trivially on the entire quotient and fixes the flag vectors.
It preserves $\Lambda'$, so it is in $N_{\Lambda',F'}$ by definition.
Both $c$ and $d$ normalize this unipotent group. Hence the
corresponding returns are in the same double-unipotent class and
give the identical adelic term. The chosen assignment is injective.
An empty fibre is finite as well. This proves the field
\texttt{fibFin} for the original map $s\mapsto s.\mathrm{toTerm}$.
\end{proof}
% END ARF1

% BEGIN ARF2
\paragraph{ARF2. Recovering an action in characteristic zero from the row index.}
Let $s_i,t_i$ ($i=1,2$) have equal arithmetic terms within each
pair $(s_i,t_i)$. Suppose $s_1,s_2$ have equal finite row indices,
and likewise $t_1,t_2$. In the decompositions of ARF1 write
$g_{t_i}=n_i m_i g_i c_i$, with $g_i=g_{s_i}$. Then
$c_2c_1^{-1}$ acts identically on $U_{F'}^\perp/U_{F'}$.
\begin{proof}
All four left lattices agree, and all four right lattices agree,
by ARF1 and the row equalities. Use their base-column transports
$a,b$ for every term. Write $L(k)$ for the tuple of reduced
middle-block matrices, computed by pairing $k\,bw_j$ with $bwd_i$
and applying each of $Y.TR$'s residue maps. QPR8 proves that on
the integral group preserving the reference triangular $(e,z)$
congruences this operation is multiplicative and has invertible
values. Indeed the span of $e$ is invariant in the reduction and
these matrices give the action on the $bw$ quotient off that span.
This works for an arbitrary typed reduction by its stated kernel
and idempotent laws. In particular $\iota_s=L(a^{-1}g_{s,f}b)$,
including terms for which that return is outside the smaller level.
The base-column transport of a rational row unipotent has $L=1$:
the strengthened congruence in PGR3 and PGR9 places its displacement
of every $bw_j$ in the $e$-span modulo the modulus. More explicitly, set
$\mu_i=g_i^{-1}m_i g_i\in N_{\Lambda',F'}$. Then
\[
 a^{-1}g_{t_i,f}b=(a^{-1}n_{i,f}a)(a^{-1}g_{i,f}b)
                        (b^{-1}\mu_{i,f}b)(b^{-1}c_{i,f}b).
\]
The first and third factors have reduced block one. Thus
\[
               \iota_{t_i}=\iota_{s_i}\,L(b^{-1}c_{i,f} b).
\]
Here $c_i$ preserves the right lattice and fixes $e'_j,z'$.
After the base-column transport it is integral and fixes these
reference vectors modulo the modulus, so all displayed matrices
are in the stated domain of $L$.

Write $c_i=v_i h_i$ in the rational splitting $N_{u'}H_{u'}$
of PGR2. Its Levi component extends the induced isometry of
$U_{F'}^\perp/U_{F'}$ by fixing the hyperbolic pairs. If
$b=h_b x_{u'}$, conjugate first by $h_b$. The integral projection
of PGR4 splits the two finite factors into lattice-preserving
factors: $h_b$ normalizes the unipotent, and $x_{u'}$ preserves
the reference lattice. Thus the unipotent has trivial reduced
block by PGR3, and the displayed $L$ is the one for $h_i$.
Cancel the two equal $\iota_{s_i}$ from the left in the equal
$\iota_{t_i}$. Multiplicativity shows that the integral element
$b^{-1}(h_2h_1^{-1})_f b$ has middle block $1$. Its $(e,z)$
congruences and the typed kernel criterion give membership in
$K_2$: in each $bw$-column the difference from the identity lies
in the $e$-span modulo the modulus, exactly as in the last step
of QPR8. Now $h_2h_1^{-1}$ is a rational element of $H_{u'}$
in $h_b K_{H,u'} h_b^{-1}$. The original column freeness PGR5
forces it to be $1$. The unipotents $v_i$ act trivially on the
quotient, which proves the assertion.
\end{proof}
% END ARF2

% BEGIN ARF3
\paragraph{ARF3. The corank and scalar test descend together.}
At the four terms of ARF2, both outputs required by
\texttt{FinX.pairIndep} agree. They are the integer $s_0$ and
the truth value of the source's scalar test.
\begin{proof}
Write $U$ for the left isotropic span and $V_i=g_i U_{F'}$.
Denote the distinguished vectors by $z$ on the left and $z_i$
in $G_i=g_iF'$. The new spans and vectors are
$V_i'=n_i V_i$, $z_i'=n_i m_i z_i$.
In particular $z_i'-n_i z_i\in V_i'$, since
$z_i\in V_i^\perp$. Similarly $n_i z-z\in U$.
Set $\beta=g_2g_1^{-1}$ and
$\tau=g_1(c_2c_1^{-1})g_1^{-1}$. Then
\[
 \beta'=g_{t_2}g_{t_1}^{-1}
      =n_2m_2\beta\tau m_1^{-1}n_1^{-1}.
\]
The middle correction $\tau$ acts identically on
$V_1^\perp/V_1$ by ARF2. For $y\in V_1^\perp$, the four maps
$m_1^{-1},\tau,\beta,m_2$ induce in succession the identity,
identity, $\bar\beta$, identity on the indicated orthogonal
quotients. Thus
$\beta'n_1 y-n_2\beta y\in V_2'$. For $x_i\in U^\perp$
one also has $n_i x_i-x_i\in U$.

To apply these congruences to the actual test spaces, put
$B_i=U+Ez+V_i+Ez_i$. This is the source's $\mathrm{UUE}$:
intersecting the two rational isotropic spans with the full row
lattice before taking their $E$-span gives those same spans by
clearing denominators. The spans for the new pair are $B_i'=n_iB_i$,
as follows from the two vector congruences above. In particular
the orthogonals $B_i^\perp$ transform by $n_i$. The paired subspace
defining $s_0$ consists of the first coordinates of pairs satisfying
\[
 x_i\in B_i^\perp,\qquad x_1-x_2\in U,\qquad
                   \beta x_1-x_2\in V_2.
\]
The bijection $(x_1,x_2)\mapsto(n_1x_1,n_2x_2)$ transports all
three conditions in both directions by the congruences. The first
projections therefore have the same dimension. Subtracting it from
the prescribed $k$ gives the same $s_0$, including when that
natural-number subtraction is zero.

The scalar test uses $y_i\in B_i\cap(U+Ez)^\perp$ and
$y_i'\in B_i\cap(V_i+Ez_i)^\perp$. Their primed-setting spaces
are again the images by $n_i$. Its two compatibility hypotheses
are $y_1-y_2\in U$ and $y_2'-\beta y_1'\in V_2$.
Both are preserved and reflected by the same two congruences.
For every pair of scalars $a,b$, the relation tested on side $i$ is
\[
            y_i-y_i'\ \in\ az+bz_i+U+V_i .
\]
Applying $n_i$ gives exactly the relation for the new distinguished
vectors, with the very same $a,b$, because the differences in $z,z_i$
lie in the displayed spans. The equivalence between the two side
relations is therefore unchanged for every admissible quadruple
and all $a,b$. This is the full scalar clause, including the
degenerate pairs. Equality of rows ensures that the original
definition of $s_0$ uses its matching-lattices branch.
\end{proof}
% END ARF3

% BEGIN ARF4
\paragraph{ARF4. The real Gram parameter and the common signature scale.}
All arithmetic terms have $G\in F_{\rm real}$, the given totally
real subfield of $E$. For every signature $\epsilon$ on $n$
coordinates there is one positive scale for the archimedean
measure valid for every ordering of the positive and negative
coordinates. These are the fields \texttt{Greal} and \texttt{cdEx}.
\begin{proof}
If $D_j$ is the leading $j$-minor of the cross-pairing matrix of
the two framed configurations, their Gram parameter is
$D_{r+1}\overline{D_{r+1}}-TT'D_r\overline{D_r}$. Both $T,T'$
are the prescribed natural numbers. Conjugation fixes this difference
and its fixed subfield is exactly the original totally real field.

Count the positive and negative coordinates, $P$ and $Q$. Choose
a sign-compatible ordering and push the given Haar measure
$D.cd(P,Q)$ through the corresponding unitary matrix reindexing
onto $U(\epsilon)$. This continuous group isomorphism produces
a nonzero Haar measure, so Haar uniqueness gives a finite positive
multiple $c$ of the specified $\mu_G(\epsilon)$. A different
ordering has the same $P,Q$ and differs by a permutation
preserving signs. Reindexing by that permutation is an automorphism
of finite order. Its pushforward multiplies Haar measure by a
positive number whose finite power is one; the number is one.
Thus the same $c$ works for every ordering. On each measurable
set the pushforward is exactly the measure of its inverse image
under the matrix reindexing in \texttt{CdTieAt}. This includes
empty positive or negative blocks and the zero-dimensional group.
\end{proof}
% END ARF4
